\section{Methods}\label{sec:methods}

\paragraph{Features.}
Each origin carries 24 base features plus one horizon-specific seasonal feature, all computed from reports dated $\le t$ (Table~\ref{tab:features}).
Missing numeric values are set to zero; reporting counts and seasonal availability indicators provide partial missingness context. The commodity vocabulary is fitted on training rows.
The cross-market features use other markets' reports of the \emph{same day}, which assumes those are published by the end of day $t$.
The older historical study examined previous-day cross-market features; that ablation is not rerun in this bounded experiment.

\begin{table*}[!t]\centering\small
\caption{Feature groups ($\ell=\ln P$); all frames end at the origin.}\label{tab:features}
\begin{tabular}{@{}lp{0.72\linewidth}@{}}\toprule
Group & Features\\\midrule
Level and lags & $\ell_t$; differences to the 1st, 2nd and 5th previous reports; days since the previous report\\
Rolling (7/30/90 d) & deviation from the rolling mean; standard deviation (30, 90); report counts (30, 90); distance to the 30-day max and min\\
Range & $(\max-\min)/\text{modal}$ on day $t$\\
Seasonal & deviation from the level at $t-365\pm3$; availability indicator; last year's change from $t-365$ to $t+h-365$\\
Cross-market & same-day mean 30-day deviation and last return of the commodity, nationally and within the state; number of reporting markets\\
Calendar & day of year (sine, cosine); day of week\\
Categorical & commodity\\\bottomrule
\end{tabular}\end{table*}

\paragraph{Baselines and learner.}
The five baselines are persistence, 7/30-day geometric rolling averages, seasonal naive with persistence fallback, and seasonal-adjusted persistence using last year's horizon change. They are evaluated on the same eligible rows as the candidate. SES, MLlib and three-seed LightGBM comparisons belong to the older historical study and are not rerun or claimed independently confirmed here.
Each horizon/variant trains a LightGBM L1 log-return model with learning rate 0.05, minimum leaf size 200, feature fraction 0.8 and bagging fraction 0.7. Early stopping uses selection-period log-return L1 loss; leaf count is selected by rupee MAE on that period. One seed and a deliberately small search budget bound this study; they do not establish best possible predictive performance or training stability.

\paragraph{Price intervals.}
Only calibration-period residuals and 30-day log-price volatility define the five volatility bins and residual quantiles. The 5/10/90/95\% quantiles produce nominal 80/90\% price intervals. This implements a baseline residual calibration, not new conformal methodology. Temporal dependence and changing distributions prevent an exchangeability-based coverage guarantee \cite{gibbs2021adaptive}; coverage is measured retrospectively, not promised for future prices.
